\documentclass[10pt,a4paper]{amsart}
\usepackage{amsmath,amssymb,amsthm,graphicx}
% Retained preamble for the Stacks Project chapter snapshot in this directory.
% The chapter's own preamble.tex is a copy of the shared Stacks Project preamble; the block
% below is sliced MECHANICALLY, line for line, from it. Nothing was added, removed or
% redefined; the wrapper class, the xr-hyper/hyperref wiring and the externaldocument list
% are omitted because the assembler supplies the document class.
%
% lines taken from preamble.tex: 3, 5, 9, 11, 12, 156, 157, 158, ... (51 lines total)
%
% For commutative diagrams we use Xy-pic
\usepackage[all]{xy}

% We use 2cell for 2-commutative diagrams.
\xyoption{2cell}
\UseAllTwocells

\newenvironment{reference}{\comment}{\endcomment}
\newenvironment{slogan}{\comment}{\endcomment}
\newenvironment{history}{\comment}{\endcomment}
\theoremstyle{plain}
\newtheorem{theorem}[subsection]{Theorem}
\newtheorem{proposition}[subsection]{Proposition}
\newtheorem{lemma}[subsection]{Lemma}
\theoremstyle{definition}
\newtheorem{definition}[subsection]{Definition}
\newtheorem{example}[subsection]{Example}
\newtheorem{exercise}[subsection]{Exercise}
\newtheorem{situation}[subsection]{Situation}
\theoremstyle{remark}
\newtheorem{remark}[subsection]{Remark}
\newtheorem{remarks}[subsection]{Remarks}
\numberwithin{equation}{subsection}
\def\lim{\mathop{\mathrm{lim}}\nolimits}
\def\colim{\mathop{\mathrm{colim}}\nolimits}
\def\Spec{\mathop{\mathrm{Spec}}}
\def\Hom{\mathop{\mathrm{Hom}}\nolimits}
\def\Ext{\mathop{\mathrm{Ext}}\nolimits}
\def\SheafHom{\mathop{\mathcal{H}\!\mathit{om}}\nolimits}
\def\SheafExt{\mathop{\mathcal{E}\!\mathit{xt}}\nolimits}
\def\Sch{\mathit{Sch}}
\def\Mor{\mathop{\mathrm{Mor}}\nolimits}
\def\Ob{\mathop{\mathrm{Ob}}\nolimits}
\def\Sh{\mathop{\mathit{Sh}}\nolimits}
\def\NL{\mathop{N\!L}\nolimits}
\def\CH{\mathop{\mathrm{CH}}\nolimits}
\def\proetale{{pro\text{-}\acute{e}tale}}
\def\etale{{\acute{e}tale}}
\def\QCoh{\mathit{QCoh}}
\def\Ker{\mathop{\mathrm{Ker}}}
\def\Im{\mathop{\mathrm{Im}}}
\def\Coker{\mathop{\mathrm{Coker}}}
\def\Coim{\mathop{\mathrm{Coim}}}
\DeclareMathSymbol{\boxtimes}{\mathbin}{AMSa}{"02}
\def\QCohstack{\mathcal{QC}\!\mathit{oh}}
\def\Cohstack{\mathcal{C}\!\mathit{oh}}
\def\Spacesstack{\mathcal{S}\!\mathit{paces}}
\def\Quotfunctor{\mathrm{Quot}}
\def\Hilbfunctor{\mathrm{Hilb}}
\def\Curvesstack{\mathcal{C}\!\mathit{urves}}
\def\Polarizedstack{\mathcal{P}\!\mathit{olarized}}
\def\Complexesstack{\mathcal{C}\!\mathit{omplexes}}
\def\Pic{\mathop{\mathrm{Pic}}\nolimits}
\def\Picardstack{\mathcal{P}\!\mathit{ic}}
\def\Picardfunctor{\mathrm{Pic}}
\def\Deformationcategory{\mathcal{D}\!\mathit{ef}}
\usepackage{cleveref}
\crefname{originalreference}{original reference}{original references}
\Crefname{originalreference}{Original reference}{Original references}
\makeatletter
\newlabel{stacks-sheaves-lemma-comparison}{{075D}{source}{}{sourcepointer-1}{}}
\newlabel{stacks-sheaves-lemma-comparison@cref}{{[originalreference][0][]075D}{[1][0][]source}{}{}{}}
\makeatother
\crefname{originalreference}{original reference}{original references}
\Crefname{originalreference}{Original reference}{Original references}
\makeatletter
\newlabel{homology-lemma-characterize-weak-serre-subcategory}{{0754}{source}{}{sourcepointer-2}{}}
\newlabel{homology-lemma-characterize-weak-serre-subcategory@cref}{{[originalreference][0][]0754}{[1][0][]source}{}{}{}}
\makeatother
\crefname{originalreference}{original reference}{original references}
\Crefname{originalreference}{Original reference}{Original references}
\makeatletter
\newlabel{stacks-morphisms-lemma-composition-smooth}{{075V}{source}{}{sourcepointer-3}{}}
\newlabel{stacks-morphisms-lemma-composition-smooth@cref}{{[originalreference][0][]075V}{[1][0][]source}{}{}{}}
\makeatother
\crefname{originalreference}{original reference}{original references}
\Crefname{originalreference}{Original reference}{Original references}
\makeatletter
\newlabel{stacks-sheaves-lemma-surjective-flat-locally-finite-presentation}{{06XH}{source}{}{sourcepointer-4}{}}
\newlabel{stacks-sheaves-lemma-surjective-flat-locally-finite-presentation@cref}{{[originalreference][0][]06XH}{[1][0][]source}{}{}{}}
\makeatother
\crefname{originalreference}{original reference}{original references}
\Crefname{originalreference}{Original reference}{Original references}
\makeatletter
\newlabel{sites-modules-lemma-tensor-product-pullback}{{03EL}{source}{}{sourcepointer-5}{}}
\newlabel{sites-modules-lemma-tensor-product-pullback@cref}{{[originalreference][0][]03EL}{[1][0][]source}{}{}{}}
\makeatother
\crefname{originalreference}{original reference}{original references}
\Crefname{originalreference}{Original reference}{Original references}
\makeatletter
\newlabel{stacks-sheaves-lemma-quasi-coherent}{{06WK}{source}{}{sourcepointer-6}{}}
\newlabel{stacks-sheaves-lemma-quasi-coherent@cref}{{[originalreference][0][]06WK}{[1][0][]source}{}{}{}}
\makeatother
\crefname{originalreference}{original reference}{original references}
\Crefname{originalreference}{Original reference}{Original references}
\makeatletter
\newlabel{sites-modules-lemma-pullback-internal-hom}{{0GN2}{source}{}{sourcepointer-7}{}}
\newlabel{sites-modules-lemma-pullback-internal-hom@cref}{{[originalreference][0][]0GN2}{[1][0][]source}{}{}{}}
\makeatother
\title{Stability properties of modules with the flat base change property}
\author{The Stacks Project Authors}
\date{stacks-project:stacks-cohomology.tex:local-snapshot-20261012; GNU FDL 1.2 or later}
\begin{document}
\maketitle
\noindent\textit{Editorial note.} All delivered mathematical text is the author's own, copied verbatim from the pinned Stacks Project chapter file stacks-cohomology.tex; no statement or proof step was rewritten, reordered or removed.\par\medskip
\noindent\textit{Editorial note.} The author marks the tensor-pullback verification in part (5) 'details omitted' with the standard Modules on Sites citation; all other parts are written out.\par\medskip
\noindent\textit{Allowed standard prerequisites (editorial).}\begin{itemize}
\item Restriction comparison on stacks. Sheaves on Stacks, Lemma stacks-sheaves-lemma-comparison. Parts (1), (3), (6).
\item Weak Serre characterization and internal hom pullback. Homology, Lemma homology-lemma-characterize-weak-serre-subcategory; Modules on Sites, Lemmas sites-modules-lemma-tensor-product-pullback and sites-modules-lemma-pullback-internal-hom. Parts (2), (5), (6).
\item Etale lifting over smooth covers. Sheaves on Stacks, Lemma stacks-sheaves-lemma-surjective-flat-locally-finite-presentation. Part (3).
\end{itemize}\medskip
\ifdefined\hypertarget\hypertarget{sourcepointer-1}{}\fi
\noindent\textit{Source reference.} Original reference 075D: sources/stacks-chow/tags.txt, source lines 9127--9127. Cross-chapter standard prerequisite cited by the author inside the counted span; it is used as a pointer and its own proof is not part of this item.\par\medskip
\ifdefined\hypertarget\hypertarget{sourcepointer-2}{}\fi
\noindent\textit{Source reference.} Original reference 0754: sources/stacks-chow/tags.txt, source lines 9118--9118. Cross-chapter standard prerequisite cited by the author inside the counted span; it is used as a pointer and its own proof is not part of this item.\par\medskip
\ifdefined\hypertarget\hypertarget{sourcepointer-3}{}\fi
\noindent\textit{Source reference.} Original reference 075V: sources/stacks-chow/tags.txt, source lines 9144--9144. Cross-chapter standard prerequisite cited by the author inside the counted span; it is used as a pointer and its own proof is not part of this item.\par\medskip
\ifdefined\hypertarget\hypertarget{sourcepointer-4}{}\fi
\noindent\textit{Source reference.} Original reference 06XH: sources/stacks-chow/tags.txt, source lines 8851--8851. Cross-chapter standard prerequisite cited by the author inside the counted span; it is used as a pointer and its own proof is not part of this item.\par\medskip
\ifdefined\hypertarget\hypertarget{sourcepointer-5}{}\fi
\noindent\textit{Source reference.} Original reference 03EL: sources/stacks-chow/tags.txt, source lines 4460--4460. Cross-chapter standard prerequisite cited by the author inside the counted span; it is used as a pointer and its own proof is not part of this item.\par\medskip
\ifdefined\hypertarget\hypertarget{sourcepointer-6}{}\fi
\noindent\textit{Source reference.} Original reference 06WK: sources/stacks-chow/tags.txt, source lines 8819--8819. Cross-chapter standard prerequisite cited by the author inside the counted span; it is used as a pointer and its own proof is not part of this item.\par\medskip
\ifdefined\hypertarget\hypertarget{sourcepointer-7}{}\fi
\noindent\textit{Source reference.} Original reference 0GN2: sources/stacks-chow/tags.txt, source lines 20777--20777. Cross-chapter standard prerequisite cited by the author inside the counted span; it is used as a pointer and its own proof is not part of this item.\par\medskip
\begin{lemma}
\label{lemma-check-flat-comparison-on-etale-covering}
Let $\mathcal{X}$ be an algebraic stack. Let $\mathcal{F}$
be an $\mathcal{O}_\mathcal{X}$-module on $\mathcal{X}_\etale$.
\begin{enumerate}
\item If $\mathcal{F}$ has the flat base change property then for any morphism
$g : \mathcal{Y} \to \mathcal{X}$ of algebraic stacks, the
pullback $g^*\mathcal{F}$ does too.
\item The full subcategory of
$\textit{Mod}(\mathcal{X}_\etale, \mathcal{O}_\mathcal{X})$
consisting of modules with the flat base change property
is a weak Serre subcategory.
\item  Let $f_i : \mathcal{X}_i \to \mathcal{X}$ be a family of
smooth morphisms of algebraic stacks such that
$|\mathcal{X}| = \bigcup_i |f_i|(|\mathcal{X}_i|)$. If each
$f_i^*\mathcal{F}$ has the flat base change property then so does
$\mathcal{F}$.
\item The category of $\mathcal{O}_\mathcal{X}$-modules
on $\mathcal{X}_\etale$ with the flat base change property
has colimits and they agree with colimits in
$\textit{Mod}(\mathcal{X}_\etale, \mathcal{O}_\mathcal{X})$.
\item Given $\mathcal{F}$ and $\mathcal{G}$ in
$\textit{Mod}(\mathcal{X}_\etale, \mathcal{O}_\mathcal{X})$
with the flat base change property then the tensor product
$\mathcal{F} \otimes_{\mathcal{O}_\mathcal{X}} \mathcal{G}$
has the flat base change property.
\item Given $\mathcal{F}$ and $\mathcal{G}$ in
$\textit{Mod}(\mathcal{X}_\etale, \mathcal{O}_\mathcal{X})$
with $\mathcal{F}$ of finite presentation and $\mathcal{G}$ having
the flat base change property then the sheaf
$\SheafHom_{\mathcal{O}_\mathcal{X}}(\mathcal{F}, \mathcal{G})$
has the flat base change property.
\end{enumerate}
\end{lemma}

\begin{proof}
Let $g : \mathcal{Y} \to \mathcal{X}$ be as in (1).
Let $y$ be an object of $\mathcal{Y}$ lying over a scheme $V$. By
Sheaves on Stacks, Lemma \ref{stacks-sheaves-lemma-comparison}
we have
$(g^*\mathcal{F})|_{V_\etale} = \mathcal{F}|_{V_\etale}$.
Moreover a comparison mapping for the sheaf $g^*\mathcal{F}$ on $\mathcal{Y}$
is a special case of a comparison map for the sheaf $\mathcal{F}$ on
$\mathcal{X}$, see
Sheaves on Stacks, Lemma \ref{stacks-sheaves-lemma-comparison}.
In this way (1) is clear.

\medskip\noindent
Proof of (2). We use the characterization of weak Serre subcategories of
Homology, Lemma \ref{homology-lemma-characterize-weak-serre-subcategory}.
Kernels and cokernels of
maps between sheaves having the flat base change property
also have the flat base change property. This is clear because
$f_{small}^*$ is exact for a flat morphism of schemes and since the
restriction functors $(-)|_{U_\etale}$ are exact (because we
are working in the \'etale topology). Finally, if
$0 \to \mathcal{F}_1 \to \mathcal{F}_2 \to \mathcal{F}_3 \to 0$
is a short exact sequence of
$\textit{Mod}(\mathcal{X}_\etale, \mathcal{O}_\mathcal{X})$
and the outer two sheaves have the flat base change property then
the middle one does as well, again because of the exactness of
$f_{small}^*$ and the restriction functors (and the 5 lemma).

\medskip\noindent
Proof of (3).
Let $f_i : \mathcal{X}_i \to \mathcal{X}$ be a jointly surjective family of
smooth morphisms of algebraic stacks and assume each $f_i^*\mathcal{F}$
has the flat base change property. By part (1), the definition of
an algebraic stack, and the fact that compositions of smooth morphisms
are smooth (see
Morphisms of Stacks, Lemma \ref{stacks-morphisms-lemma-composition-smooth})
we may assume that each $\mathcal{X}_i$ is representable by a scheme.
Let $\varphi : x \to x'$ be a morphism of $\mathcal{X}$ lying over
a flat morphism $a : U \to U'$ of schemes. By
Sheaves on Stacks, Lemma
\ref{stacks-sheaves-lemma-surjective-flat-locally-finite-presentation}
there exists a jointly surjective family of \'etale morphisms
$U'_i \to U'$ such that $U'_i \to U' \to \mathcal{X}$ factors through
$\mathcal{X}_i$. Thus we obtain commutative diagrams
$$
\xymatrix{
U_i = U \times_{U'} U_i' \ar[r]_-{a_i} \ar[d] &
U_i' \ar[r]_{x_i'} \ar[d] & \mathcal{X}_i \ar[d]^{f_i} \\
U \ar[r]^a & U' \ar[r]^{x'} & \mathcal{X}
}
$$
Note that each $a_i$ is a flat morphism of schemes as a base change of $a$.
Denote $\psi_i : x_i \to x'_i$ the morphism of $\mathcal{X}_i$ lying over
$a_i$ with target $x_i'$. By assumption the comparison maps
$c_{\psi_i} :
(a_i)_{small}^*\big(f_i^*\mathcal{F}|_{(U'_i)_\etale}\big)
\to f_i^*\mathcal{F}|_{(U_i)_\etale}$ is an isomorphism.
Because the vertical arrows $U_i' \to U'$ and $U_i \to U$ are \'etale,
the sheaves $f_i^*\mathcal{F}|_{(U_i')_\etale}$ and
$f_i^*\mathcal{F}|_{(U_i)_\etale}$ are the restrictions of
$\mathcal{F}|_{U'_\etale}$ and $\mathcal{F}|_{U_\etale}$
and the map $c_{\psi_i}$ is the restriction of $c_\varphi$ to
$(U_i)_\etale$, see
Sheaves on Stacks, Lemma \ref{stacks-sheaves-lemma-comparison}.
Since $\{U_i \to U\}$ is an \'etale covering, this implies
that the comparison map $c_\varphi$ is an isomorphism which is what
we wanted to prove.

\medskip\noindent
Proof of (4). Let $\mathcal{I} \to
\textit{Mod}(\mathcal{X}_\etale, \mathcal{O}_\mathcal{X})$,
$i \mapsto \mathcal{F}_i$ be a diagram and assume each $\mathcal{F}_i$
has the flat base change property. Let $\varphi : x \to x'$ be a morphism
of $\mathcal{X}$ lying over the flat morphism of schemes $f : U \to U'$.
Recall that $\colim_i \mathcal{F}_i$ is the sheafification of the presheaf
colimit. As we are using the \'etale topology, it is clear that
$$
(\colim_i \mathcal{F}_i)|_{U_\etale} =
\colim_i {\mathcal{F}_i}|_{U_\etale}
$$
and similarly for the restriction to $U'_\etale$. Hence
\begin{align*}
f_{small}^*((\colim_i \mathcal{F}_i)|_{U'_\etale})
& =
f_{small}^*(\colim_i {\mathcal{F}_i}|_{U'_\etale}) \\
& =
\colim_i f_{small}^*({\mathcal{F}_i}|_{U'_\etale}) \\
& \xrightarrow{\colim c_\varphi}
\colim_i \mathcal{F}_i|_{U_\etale} \\
& =
(\colim_i \mathcal{F}_i)|_{U_\etale}
\end{align*}
For the second equality we used that $f_{small}^*$ commutes with colimits
(as a left adjoint). The arrow is an isomorphism as each
$\mathcal{F}_i$ has the flat base change property. Thus the colimit
has the flat base change property and (4) is true.

\medskip\noindent
Part (5) holds because tensor products commute with pullbacks, see
Modules on Sites, Lemma \ref{sites-modules-lemma-tensor-product-pullback}.
Details omitted.

\medskip\noindent
Let $\mathcal{F}$ and $\mathcal{G}$ be as in (6). Since $\mathcal{F}$
is quasi-coherent it has the flat base change property by
Sheaves on Stacks, Lemma \ref{stacks-sheaves-lemma-quasi-coherent}.
Let $\varphi : x \to x'$ be a morphism of $\mathcal{X}$ lying
over the flat morphism of schemes $f : U \to U'$.
As we are using the \'etale topology, we have
$$
\SheafHom_{\mathcal{O}_\mathcal{X}}(\mathcal{F}, \mathcal{G})|_{U_\etale} =
\SheafHom_{\mathcal{O}_U}(\mathcal{F}|_{U_\etale}, \mathcal{G}|_{U_\etale})
$$
and similarly for the restriction to $U'_\etale$ (details omitted). Hence
\begin{align*}
f_{small}^*(
\SheafHom_{\mathcal{O}_\mathcal{X}}(\mathcal{F}, \mathcal{G})|_{U'_\etale})
& =
f_{small}^*(
\SheafHom_{\mathcal{O}_{U'}}(
\mathcal{F}|_{U'_\etale}, \mathcal{G}|_{U'_\etale})) \\
& =
\SheafHom_{\mathcal{O}_{U'}}(
f_{small}^*(\mathcal{F}|_{U'_\etale}),
f_{small}^*(\mathcal{G}|_{U'_\etale})) \\
& \xrightarrow{c_\varphi}
\SheafHom_{\mathcal{O}_U}(\mathcal{F}|_{U_\etale}, \mathcal{G}|_{U_\etale}) \\
& =
\SheafHom_{\mathcal{O}_\mathcal{X}}(\mathcal{F}, \mathcal{G})|_{U_\etale}
\end{align*}
Here the second equality is
Modules on Sites, Lemma \ref{sites-modules-lemma-pullback-internal-hom}
which uses that $f : U \to U'$ is flat and hence the morphism
of ringed sites $f_{small}$ is flat too.
The arrow is an isomorphism as both $\mathcal{F}$ and $\mathcal{G}$
have the flat base change property. Thus our $\SheafHom$
has the flat base change property too as desired.
\end{proof}

\end{document}
